\section{The Sonnet Challenge}

As a condition for initiation into the Fedeli d'Amore, the candidate submitted a sonnet. Writing a sonnet has several benefits for the poet:

\begin{itemize}


\item He cannot rely on discursive thinking to express ideas. Instead, he has to create images and thoughts in the reader's mind through succinct sentences.

\item He is forced to write with a high idea density. Since the number of words is limited, there can be no unnecessary words.

\item The restriction to a rhyme scheme and meter disciplines the mind.
\end{itemize}


Fundamentally, some thoughts can only be expressed in a poem.

In that spirit, several participants in our esoteric group have accepted the sonnet challenge. We shall be posting them each week for the next month or so.


\flrightit{Posted on 2023-01-18 by Cologero}

\begin{center}* * *\end{center}

\begin{footnotesize}
\begin{sffamily}
\texttt{Ignatius on 2023-01-19 at 09:21 said:}

\begin{verse}
Die Zeit hat nur ihn selbst zurückgelassen;\\ nie wird er den Sphären lauschen wie an jenen Tagen;\\ wenn er doch nie die höheren Mysterien aufgeschlagen;\\ wär ihm das bescheidne Glück belassen.

Welch Grausamkeit, sein Leben so zu schließen;\\ er verpaßte Anfang, Mitte und End in einem;\\ weder konnte er in Gottes Gnade sich vereinen;\\ noch die frischen Wonnen des Tages genießen.

Das Schicksal ist die kalte Zwischenwelt;\\ Denen es bei den Menschen nicht gefällt;\\ Nicht tot, nicht lebend in die Höhe starren.

Kann sich nicht winden und nicht beugen;\\ Nur noch aus sich selbst heraus erzeugen;\\ Um auf ewig pflanzlich auszuharren.
\end{verse}

\hfill

\end{sffamily}
\end{footnotesize}

